%% kapa: exam=1 round=30 year=2019 period=1 form=A subject=civil_law subject_ko=민법 seq=30 number=30
%% source: Q-Net 감정평가사 기출문제 (cst00309) artlSeq=5205205 file=extracted/1008/5205205/2019년 제30회 감정평가사 1차시험 문제지/2019년 제30회 감정평가사 1차시험 문제지/2019년도 제30회 감정평가사 1차 1교시 A형.hwp
%% conversion: kordoc 4.24.51 -> tools/gen1.py (auto); review=proofread-sonnet; flags=none
\begin{question}{30}
점유자와 회복자의 관계에 관한 설명으로 옳지 않은 것은? (다툼이 있으면 판례에 따름)
\begin{choices}
  \item 과실을 수취한 자가 선의의 점유자로 보호되기 위해서는 과실수취권을 포함하는 권원이 있다고 오신할만한 정당한 근거가 있어야 한다.
  \item 폭력 또는 은비에 의한 점유자는 수취한 과실을 반환하여야 한다.
  \item 점유물이 점유자의 책임있는 사유로 인하여 멸실 또는 훼손한 때에는 선의의 자주점유자라도 그 손해의 전부를 배상하여야 한다.
  \item 악의의 점유자도 점유물을 반환할 때에는 회복자에 대하여 필요비의 상환을 청구할 수 있다.
  \item 선의의 점유자가 과실을 취득한 경우에는 통상의 필요비는 청구하지 못한다.
\end{choices}
\end{question}
\answer{3}
